% subsection 4.3
\subsubsection{Etapa 3: Evaluación de calidad}

De acuerdo con la literatura sobre estudios de mapeo sistemático como el de~\cite{Ali-01} y~\cite{Petersen2008}, la evaluación de calidad 
no es un requisito obligatorio en este tipo de investigaciones; sin embargo, su incorporación permite aumentar la rigurosidad del proceso 
de selección y análisis de estudios. Por esta razón, en este trabajo se realizó una evaluación de calidad con el objetivo de verificar la 
relevancia y el impacto de los estudios identificados en relación con los objetivos del SMS.\@

Para llevar a cabo esta evaluación se utilizaron los índices CVI (\textit{Content Validity Index}), SCI (\textit{Study Citation Index})
e IRRQ (\textit{Index of Relationship to Research Questions}), los cuales fueron definidos y descritos en detalle durante la fase
de planificación (Sección~\ref{subsec:criterios_calidad}). Los valores de estos índices fueron calculados con el apoyo de la herramienta
SMS-Builder~\cite{Candela-Uribe2022}.

%sub-subsection 4.3.1
\subsubsection{Evaluación de validez del contenido}

En esta etapa se analizó el contenido de los estudios seleccionados con el fin de determinar su valor dentro del contexto de la 
investigación. Para ello se aplicó el índice CVI, el cual permite evaluar la relevancia del contenido de cada estudio respecto al 
dominio de investigación analizado.

Posteriormente, se efectuó un análisis de frecuencia para identificar el cuartil de estudios más significativos, es decir, aquellos 
con los valores más altos del índice CVI.\@ La definición formal de esta métrica se presenta en la 
Sección~\ref{subsec:criterios_calidad}, específicamente en la Ecuación~\ref{eq:cvi}.

%sub-subsection 4.3.2
\subsubsection{Índice de calidad basado en el número de citas}

La evaluación del impacto científico de los estudios se realizó mediante el índice SCI.\@ Este índice considera el número de citas
recibidas por cada publicación como un indicador de su influencia dentro de la comunidad académica.

El cálculo del índice SCI se realizó utilizando SMS-Builder en conjunto con los datos de citación obtenidos a partir de
Google Académico. Posteriormente,se llevó a cabo un análisis de frecuencia orientado a identificar el cuartil de estudios con 
mayor relevancia según este indicador. La formulación formal de esta métrica se presenta en la Sección~\ref{subsec:criterios_calidad}, 
específicamente en la Ecuación~\ref{eq:sci}.

%sub-subsection 4.3.3
\subsubsection{Relación de los estudios con las preguntas de investigación}

Finalmente, se aplicó el índice IRRQ con el fin de evaluar el grado de relación entre los estudios seleccionados y las preguntas de
investigación planteadas en este estudio.

Para ello, la herramienta SMS-Builder fue configurada para establecer la relación directa entre cada estudio y los temas de clasificación
definidos en el SMS.\@ Como en los casos anteriores, se realizó un análisis de frecuencia que permitió identificar el cuartil de
estudios con mayor relación con las preguntas de investigación. La formulación formal de esta métrica se presenta en la Sección~\ref{subsec:criterios_calidad}, 
específicamente en la Ecuación~\ref{eq:irrq}.